\documentclass[10pt,a4paper,twoside,bg=print,twocolumn,openany,nodeprecatedcode]{dndbook}
\usepackage[utf8]{inputenc}
\usepackage{tabulary}
\usepackage[normalem]{ulem}
\usepackage{color,soul}
\usepackage{float}
\usepackage{caption}
\usepackage{wrapfig}
\usepackage{tocloft}
\usepackage[toc]{multitoc}
\usepackage[hidelinks]{hyperref}
\raggedbottom
\extrafloats{2000}
\cftsetindents{section}{0em}{0em}
\cftsetindents{subsection}{0em}{0em}
\cftsetindents{subsubsection}{0.2em}{0em}
\setcounter{tocdepth}{1}
\renewcommand*{\multicolumntoc}{3}
\setlist[description]{nolistsep,listparindent=3pt,topsep=6pt}
\date{}
\begin{document}
\frontmatter
\tableofcontents
\mainmatter
\chapter{Descent into the Lost Caverns of Tsojcanth}
\DndDropCapLine{D}{}eep in the Yatil Mountains lie the Lost Caverns of Tsojcanth, formerly occupied by the legendary archmage Iggwilv the Witch Queen. Though Iggwilv is long gone, her lair is anything but empty. Demons, giants, and other formidable creatures haunt the perilous caverns, and the archmage's magical defenses remain intact. The rewards for braving these threats defy imagination. Iggwilv is rumored to have amassed a magical hoard of unsurpassed value, a trove of such fame that scores of adventurers have perished in search of it.


This supplement is part of a yearlong celebration of Dungeons \& Dragons and its 50th anniversary. The adventure presented herein is an abridged, quick-play version of "The Lost Caverns of Tsojcanth," which appears in Quests from the Infinite Staircase. See that book for the full adventure.

"Descent into the Lost Caverns of Tsojcanth" is designed for four to six 9th-level characters.

\section{Using This Supplement}
In addition to the adventure, this supplement includes the following appendices:

Appendix A presents a small trove of magic items available to all who brave the Lost Caverns of Tsojcanth. At the DM's discretion, the characters might earn one or more of these treasures as a reward for completing the adventure.

Appendix B contains a creature that dwells in the Lost Caverns, the \textbf{pech}.

Appendix C features a score sheet for use in competitive play. See the "Tournament Rules" section for how to run the adventure as a tournament-style scenario.

Appendix D provides six premade characters suitable for this adventure. Each player can choose one of these characters or provide a 9th-level character of their own.

\begin{DndReadAloud}{About the Original}
Published in 1982, the official \textit{Lost Caverns of Tsojcanth} module was an expanded and revised version of \textit{The Lost Caverns of Tsojconth}, a tournament adventure created by Gary Gygax for Winter Con V in 1976. The adventure offers tantalizing details about Iggwilv the Witch Queen—perhaps better known as the archmage Tasha, famous for her \textit{Hideous Laughter} spell. This adaptation pays homage to the original by including a score sheet for competitive play.

\end{DndReadAloud}

\subsection{Running the Adventure}
To run the adventure, you need the fifth edition core rulebooks: the \textit{Player's Handbook}, \textit{Dungeon Master's Guide}, and \textit{Monster Manual}. When a creature's name appears in \textbf{bold} type, that's a visual cue pointing you to its stat block as a way of saying, "Hey, DM, you should get this creature's stat block ready. You're going to need it." The \textit{Monster Manual} contains stat blocks for most of the creatures in these adventures; the \textbf{pech} appears in appendix B.

\begin{DndReadAloud}{}
Text that appears in a box like this is meant to be read aloud or paraphrased for the players when their characters first arrive at a location or under a specific circumstance, as described in the text.

\end{DndReadAloud}

Spells and equipment mentioned in this supplement are described in the \textit{Player's Handbook}. Magic items are described in the \textit{Dungeon Master's Guide}.

\section{Background}
Iggwilv carved out a kingdom for herself through magical prowess, strength, and wit, bolstered by a horde of demons bound by magic to further her aims. In her quest for absolute power, Iggwilv accumulated countless enemies determined to destroy her. Among these was the demon lord Graz'zt, whom the archmage imprisoned for a time. Eventually, Iggwilv's ambitions proved excessive, and a climactic battle with Graz'zt broke her power. The Witch Queen then vanished. As her evil influence waned, Iggwilv's realm was sundered, and her lair was lost.

This lost lair purportedly lies deep in a series of caverns somewhere in the Yatil Mountains. In the decades since Iggwilv's disappearance, adventurers and treasure hunters have sought the Lost Caverns of Tsojcanth and the treasures within. This adventure assumes the characters know the way to the fabled lair, having assembled fragments of information regarding the caverns' location or learned it from another reliable source.

\subsection{Adventure Hooks}
If the prospect of claiming Iggwilv's hoard isn't enough to persuade the characters to undertake this quest, you can further motivate the party with either of the following adventure hooks:


\begin{description}
\item[Iggwilv's Rival.]
An archmage who was once embroiled in a heated rivalry with Iggwilv has reason to believe she isn't dead. The archmage recruits the characters to search the Lost Caverns for any evidence of the Witch Queen's lingering presence.
\item[Margrave's Agents.]
The margrave of a nation near the mountains has uncovered clues to the Lost Caverns' whereabouts, thanks to a few partially successful expeditions in recent years. Most of the agents the margrave sent came back battle worn and empty-handed, if they returned at all. The margrave hires the characters to recover Iggwilv's hoard before it falls into the hands of his rivals.
\end{description}

\subsection{Setting the Adventure}
On the world of Oerth in the Greyhawk campaign setting, the Lost Caverns hide within the Yatil Mountains, nestled between the nations of Ket and Perrenland, with Bissel to the south.

You can set the adventure in a different world by placing it in any mountain range near a kingdom or settlement—before her disappearance, Iggwilv's influence spanned multiple worlds and planes of existence. Consider the following suggestions:


\begin{description}
\item[Dragonlance.]
The Speaker of the Suns, leader of the Qualinesti elves, might charge the characters to search for the Lost Caverns in the Kharolis Mountains south of the elves' forest home.
\item[Eberron.]
The nations of Aundair, Breland, and Thrane vie to explore the Lost Caverns in the Blackcap Mountains. Iggwilv's experiments there drew power from the depths of Kyber.
\item[Forgotten Realms.]
The Lost Caverns are hidden in the Sunset Mountains, northeast of Proskur along the High Road. The High Overseer of Elturel might hire the characters to locate the caverns.
\end{description}

\section{Tournament Rules}
If you wish to emulate the competitive adventures of old, you can run this adventure as a tournament-style scenario by following the guidance in this section. Rather than pitting characters against one another, adventuring parties are scored as a group.

These rules are optional. You don't need to score your players to run this adventure. The scoring system included in this supplement is meant to spark joy, not contention.

\subsection{Rewards}
At game conventions and store events, these rules can be used to award prizes to the highest-scoring groups. Otherwise, they mainly exist for groups to enjoy bragging rights, comradery, and other intangible benefits. Consider sharing your group's final score as a lighthearted challenge to other players.

\subsection{Scoring}
Have this adventure's score sheet (see appendix C) handy as you run the adventure. As the characters traverse the Lost Caverns, silently cross-reference the actions of the party with those denoted on the score sheet, making note of any events that qualify for a point increase or deduction.

Keep these actions and their associated point values secret during the adventure. You can share them with the players afterward, along with the group's final score. Alternatively, you can announce these point fluctuations as they occur to your players for comedic or dramatic effect.

After the adventure, use the score sheet to tally the group's points and calculate their final score, then report that score to the players. It's possible—albeit unlikely—to earn a negative final score.

\subsection{Time Limit}
As a tournament scenario, this adventure has a strict time limit: a single game session equal to four hours of play. Encourage players at the onset not to dawdle—the clock is ticking! The adventure concludes once that time has elapsed or whenever any of the following conditions have been met:


\begin{itemize}
\item Every character in the party is dead or otherwise out of commission.
\item The characters explore every area of the lesser caverns (presented later in this adventure).
\item The characters descend to the greater caverns (see area L16 of the lesser caverns) or leave the Lost Caverns altogether. Let players know that these actions will end the adventure (and thus their opportunity to score more points).
\end{itemize}

\subsubsection{Resting}
To keep the dungeon challenging, the characters can take only one short rest during the session. Long rests are prohibited for tournament play of this adventure.

\section{Lost Caverns}
The Lost Caverns of Tsojcanth hide in the Yatil Mountains south of Iggwilv's Horn. Not much is known about the caverns' namesake, the mysterious Tsojcanth, who was eclipsed by the archmage who succeeded them. Whether Tsojcanth was slain by Iggwilv or left before her arrival, this predecessor's only remaining trace is the caverns' name, while the Witch Queen's influence persists long after her departure.

The caverns are divided into the lesser and greater caverns. This abbreviated adventure focuses only on the former.

\subsection{Starting the Adventure}
Read or paraphrase the following text to begin the adventure:

\begin{DndReadAloud}{}
You are a member of a group of adventurers, united in weeks past by your pursuit of a common prize: the treasures of the Lost Caverns of Tsojcanth—the former lair of the archmage Iggwilv, the so-called Witch Queen.

Guided by a series of hard-won clues, your trek through the Yatil Mountains halts at a yawning cleft in the mountainside. Stalactites and stalagmites growing at the entrance give the cave mouth the impression of fangs waiting to snap together. The opening is forty feet wide, twenty feet high, and blackened by soot. Past the stone maw, a wide passage with roughly hewn stairs descends into darkness.

\end{DndReadAloud}

The tunnel slants down to area L1 of the lesser caverns.

\subsection{Lesser Caverns Features}
The lesser caverns form the upper level of the Lost Caverns. Unless otherwise noted, the lesser caverns have the following features:


\begin{description}
\item[Ambience.]
The caverns are damp, humid, and full of life. Stalactites drip into shallow puddles, and water trickles down the slick cave walls. The sounds of cave bats, small insects, and the caverns' monstrous denizens echo throughout.
\item[Ceilings.]
Ceilings are 15 feet high in corridors and tunnels and 20 feet high in chambers.
\item[Lighting.]
The caverns are unlit. Denizens carry their own light or rely on darkvision to see. Area descriptions assume the characters have a light source or other means of seeing in the dark.
\item[Underground Rivers.]
The rivers that span the lesser caverns have strong currents. A creature that enters a river for the first time on a turn or starts its turn there must succeed on a DC 15 Strength saving throw or be swept 20 feet in the direction indicated by the arrows on the map.
\end{description}

\subsection{Lesser Cavern Locations}
The following locations are keyed to map 1.1.

\subsubsection{L1: Entry Caverns}
\begin{DndReadAloud}{}
This wide chamber is ringed with a mixture of rough-hewn and natural passageways leading away from the entrance and into the unknown. Beside each of these tunnels is a grotesque face carved into the stone. Though each face is slightly different—one has flap-like ears, another protruding tusks, and a third drooping wattles—all are strange and doleful.

\end{DndReadAloud}

Each of the carved faces is magical. They are immune to all damage and speak a few programmed messages. When a creature approaches one of the tunnels leaving the chamber, the face beside it animates and warns in a deep, dire tone, "Turn back—this is not the way!"

A character who examines one of the faces and succeeds on a DC 15 Wisdom (Perception) check notices the glint of a gemstone in its mouth.

\subparagraph{Speaking to the Faces}
The faces repeat their warning if spoken to, except in the following two circumstances:


\begin{description}
\item["Say 'Aah!]
'" If a character asks one of the faces to open its mouth, stick out its tongue, or yawn, the face sticks out its tongue and says "aah," revealing a gemstone the size of a hen's egg on its tongue.
\item["Truth.]
" If a character asks one of the faces which way to go or uses the word "truth" while speaking near a face, the faces respond in unison. Five of the faces lie, saying, "My way is the right way." The face marking the southeast tunnel leading to area L9, however, responds, "I watch the best way."
\end{description}

\subparagraph{Treasure}
Each of the faces contains a different gem worth 200 gp: amber, amethyst, blue aquamarine, garnet, peridot, and pink tourmaline.

A creature that tries to retrieve a gem from one of the faces without first asking the face to open its mouth must make a DC 17 Dexterity (Sleight of Hand) check. On a failed check, the creature takes 11 (2d10) bludgeoning damage as the mouth clamps shut. The mouth similarly damages tools or weapons inserted into it without its permission.

\subsubsection{L2: Slate Chamber}
\begin{DndReadAloud}{}
Alternating blue and gray bands of slate and shale compose the walls of this cave. Rusted weapons lie scattered on the ground.

A hulking, crudely formed figure made of gray clay charges at you from an alcove in the northeast, growling. The hilt of a pristine sword protrudes from the creature's chest.

\end{DndReadAloud}

The figure is a \textbf{clay golem} created by Iggwilv to destroy intruders. The golem attacks any creature that enters the chamber. The golem fights to the death and pursues trespassers relentlessly, even chasing them out of the caverns.

\subparagraph{Treasure}
The weapon stuck in the clay golem is a slender-bladed Giant Slayer Shortsword. The sword's name ("Longtooth") is engraved into the blade in Gnomish. Once the golem is destroyed, the sword can easily be freed from its chest.

\subsubsection{L3: Guano-Covered Cave}
\begin{DndReadAloud}{}
The floor of this vaulted cave is covered in pungent piles of bat guano. Large leather hides hang from the stalactites above, flapping in a chilly draft. Sounds of flowing water echo from the far side of the cave.

\end{DndReadAloud}

The ceiling of this soiled cave is 50 feet high. This cave once hosted dozens of bats, but a \textbf{cloaker} that dwells here devoured them. The hungry creature hides among the stalactites, waiting to fall on one of the characters.

\subparagraph{Tunnel Ledge}
A ledge on the north wall of the cave stretches out 30 feet above the floor. A short tunnel leads from the ledge to the river beyond. The water's surface is 40 feet below the tunnel's exit.

\subsubsection{L4: Littered Cave}
\begin{DndReadAloud}{}
This small cave is littered with broken bones. Near the back wall lies a busted skeleton with a horned bovine head splintered with cracks. Two leather sacks brimming with gold rest near the skeleton.

\end{DndReadAloud}

The skeleton is that of a minotaur thief who fell prey to the trap in this cave. The two leather sacks beside it are illusions designed to lure greedy intruders to their doom. A character who inspects the sacks and succeeds on a DC 18 Intelligence (Investigation) check or who makes physical contact with one of them discerns their false nature.

\subparagraph{Crushing Ceiling Trap}
A 20-foot-wide section of the ceiling over the skeleton contains a crushing deadfall. A character who examines the ceiling and succeeds on a DC 17 Wisdom (Perception) check notices bits of decayed flesh stuck to the ceiling, which is flat and made of stone.

Whenever a creature enters the area marked on the map, the trap magically triggers. That section of the ceiling rapidly drops, crushing anything in its path. Creatures in that area must make a DC 17 Dexterity saving throw. On a failed save, a creature takes 55 (10d10) bludgeoning damage and has the prone condition. On a successful save, it takes half as much damage only and is pushed to the nearest unoccupied space outside the trap's area. Objects wedged under the trap are destroyed. The stone then retracts into the ceiling, and the trap resets.

A successful \textit{Dispel Magic} spell (DC 17) suppresses the trap and its illusions for 24 hours.

\subparagraph{Treasure}
A pouch on the minotaur's skeleton holds five diamonds worth 100 gp each. The illusory sacks contain nothing.

\subsubsection{L5: Grotto}
\begin{DndReadAloud}{}
This tunnel has been worked recently, judging by the stone chips and dust covering the floor. Roughly carved images of cave-dwelling creatures decorate the smooth stone walls.

The sounds of mining tools echo from the end of the tunnel, where hewn steps rise into the darkness.

\end{DndReadAloud}

Eight pechs (see appendix B) are diligently carving a tunnel from this grotto up into the mountain toward an enclosed cavern they can sense. For now, the tunnel leads only to a dead end. If the characters draw attention with light or noise, the pechs investigate.

The pechs are initially indifferent toward the characters, admonishing them to douse any bright lights they carry, as the lights hurt the pechs' eyes. If the characters comply, the pechs become curious and willing to chat. Otherwise, they resume their work, defending themselves if necessary.

In conversation, the pechs warn the characters about the clay golem in the slate chamber (area L2). Characters can improve the pechs' attitude toward them by offering the pechs at least 500 gp of gems or a unique carving or sculpture, especially one made of stone. Friendly pechs might be willing to help fight the clay golem or use their Communal Spellcasting action to restore petrified characters.

\subparagraph{Carvings}
The carvings on the walls depict cave-dwelling creatures, such as bats and crickets, as well as creatures from the Elemental Plane of Earth such as \textbf{dao}, pechs, \textbf{xorn}, and earth elementals. Characters who inspect the carvings and succeed on a DC 16 Intelligence (Arcana) check identify the planar creatures.

\subsubsection{L6: Fungal Cavern}
\begin{DndReadAloud}{}
Bioluminescent fungi, some as high as six feet tall, fill this moist cavern. The caps of the mushrooms sprouting from the floor and walls glow with a faint blue light. Pale, dog-sized insects chitter among the fungus.

\end{DndReadAloud}

Three ghost-white cave crickets (use the \textbf{giant frog} stat block) nibble at the glowing mushrooms in this cave, which radiate dim light. The crickets are indifferent toward the characters.

If the characters disturb the crickets, the insects begin to chirp loudly. The trolls in area L8 investigate any loud noises coming from this area.

\subparagraph{Mushrooms}
The fungi in this chamber are edible, and characters can forage here without having to make an ability check. Doing so doesn't disturb the crickets.

\subsubsection{L7: Slimy Cavern}
\begin{DndReadAloud}{}
This cavern is full of sizable mushrooms, but the air stings with an acrid scent. Slime in sickly yellow, green, and orange shades blankets most of the fungi, as well as the walls and ceiling.

Near the back of the cave is a cranny mostly free of fungus. A black-cloaked figure, motionless, is wedged into the crack.

\end{DndReadAloud}

Originally this was another cultivated fungus farm like area L6, but slime molds are overtaking the cave. The cloaked figure is the corpse of an elf adventurer who died here, trapped by the green slime. His body has been partially calcified by the mineral-rich water dripping from the ceiling.

\subparagraph{Green Slime}
Six patches of green slime (see the \textit{Dungeon Master's Guide}) cling to the ceiling near the fallen adventurer. A character who inspects the slime and succeeds on a DC 15 Intelligence (Nature) check recognizes the danger it poses. The slime falls on creatures that approach the corpse.

\subparagraph{Mushrooms}
The fungi here are mostly edible. Foraging in this cavern requires a successful DC 15 Wisdom (Survival) check to avoid the harmful molds and slime. On a failed check, the mushrooms the character gathers are contaminated by the slime. A creature that eats a contaminated mushroom must succeed on a DC 12 Constitution saving throw or have the poisoned condition for 8 hours.

\subparagraph{Treasure}
The elf corpse carries a silvered \textit{rapier} and wears \textit{Bracers of Defense}.

\subsubsection{L8: Stinking Cave}
\begin{DndReadAloud}{}
This reeking cave is filled with rotting vegetation, bones, pieces of pale chitin, dung, and less identifiable foulness. Three larger heaps of material serve as giant sleeping pallets.

Three hulking creatures with warty green skin and wiry hair mill through the detritus, chewing cracked bones clutched in their clawed hands.

\end{DndReadAloud}

This cave is the lair of three trolls. They regularly gather rotting material from their den and spread it through the nearby fungus caverns (areas L6 and L7) to ensure the mushrooms grow and attract cave crickets for the trolls to feed on. The trolls are ornery, having long gone without the taste of humanoid flesh—their favorite delicacy. The trolls attack anyone they notice, fantasizing aloud in Giant about the meal they're about to have.

\subparagraph{Treasure}
Gathered under the trolls' beds is their accumulated loot: 120 gp, 900 sp, four gems worth 50 gp each, gnome-sized wooden dentures set with mother-of-pearl teeth (worth 50 gp), a \textit{quiver} with twenty silvered Arrows (20), and a Potion of Greater Healing.

\subsubsection{L9: Bat Corridor}
\begin{DndReadAloud}{}
The sandy floor of this long, wide corridor is strewn with piles of bat guano. Small fungi sprout from these deposits and the walls. A cacophony of chitters and squeaks echoes from above.

\end{DndReadAloud}

The ceiling of this corridor is covered in thousands of roosting bats. Bright light or any noise louder than a whisper agitates the bats and sends them into a panic, causing a number of swarms of bats equal to the number of characters in this area descend from the ceiling and attack the source of the disturbance. The swarms extinguish any nonmagical light sources they come within 5 feet of with a flurry of buffeting wings.

Whenever a creature destroys a swarm, a new swarm immediately forms in an unoccupied space within 15 feet of the creature. The new swarm acts on the destroyed swarm's initiative count.

If the characters extinguish all light sources and cease attacking and making loud noises, the bats settle back into their roosts at the end of the round. They don't pursue creatures beyond the corridor.

\subsubsection{L10: Long Gallery}
\begin{DndReadAloud}{}
The cavern widens into a gallery of considerable height and length. Clusters of small fungi grow throughout the chamber, and odd indentations dot the floor and walls at semiregular intervals. Somewhere farther back in the cave, a dim glow casts flickering shadows.

\end{DndReadAloud}

The indentations are the rock-cyst burrows of four cave morays—sluglike scavengers that slither between the burrow entrances through tunnels in the rock. The cave morays use the \textbf{giant poisonous snake} stat block. They use their reach to bite at prey from the safety of their burrows.

Cave morays inside their burrows have \textit{three-quarters cover}. A creature within 5 feet of a burrow can pull a cave moray from its burrow by grappling it. If removed from its burrow, a cave moray frantically tries to escape back to its hole.

Loud sounds here, including the din of battle, draw the attention of the fomorians in area L11. The cave morays retreat into their burrows when the fomorians arrive.

\subsubsection{L11: Fomorian Cave}
\begin{DndReadAloud}{}
This cave stinks of sweat, smoke, and spoiled meat. Old, cracked bones and skulls lie strewn on the floor. A large, flat rock supports massive stone goblets and the carcass of a giant toad, picked clean and buzzing with flies. Two piles of filthy hides and skins form beds at the far end of the cave.

Two hideous giants squat near a circle of hefty stones enclosing a crackling fire. Above the roaring flames, a charred giant insect rotates slowly on a spit.

\end{DndReadAloud}

Two fomorians inhabit this cave. The beds, fire pit, and utensils are all sized for Huge creatures.

The fomorians are cruel, violent, and bored. When the giants notice the characters, they abandon their meal and attack. One of the fomorians wields the roasting spit as its weapon, tearing bites off the bug-tipped skewer between swings.

\subparagraph{Treasure}
Stashed beneath the fomorians' beds are two ivory mammoth tusks (each is worth 600 gp and weighs 100 pounds), a \textit{Spell Scroll} of \textit{Levitate}, and a beaten copper bowl with lapis lazuli handles, which is worth 750 gp.

\subsubsection{L12: River Ledge}
\begin{DndReadAloud}{}
A ten-foot-long wooden boat is tied to a ledge near the water. Intricate, flowing knotwork carving decorates the sides of the craft, which tapers to a point at both ends. Two oars lie inside.

\end{DndReadAloud}

This ledge is visible from the lake (area L13) and from the shores of areas L9 and L14.

\subparagraph{Treasure}
The vessel is a \textit{Folding Boat} created to traverse the lake and the rivers that branch from it. In addition to its usual properties, the boat has an extra command word that functions only while the boat is within the Lost Caverns. A character who knows this additional command word ("Shrimpkin")—discovered by casting the \textit{Identify} spell on the boat or by examining clues in the boat's intricate designs and succeeding on a DC 16 Intelligence (Arcana) check—can command the boat to row itself in a particular direction at a speed of 30 feet or to hold its position in defiance of the water's swiftest currents.

\subsubsection{L13: Underground Lake}
\begin{DndReadAloud}{}
A river flows into the south side of this high-vaulted cavern, feeding a glassy, ebon-hued lake. Drops of water fall from the ceiling high above, rippling the mirrorlike surface of the water.

Two rivers flow out of the northern side of the lake, while a third flows out to the west. Ledges along the shore lead to dry tunnels.

\end{DndReadAloud}

The lake teems with blind cave fish and crayfish. These critters are a food source for many of the denizens in the lesser caverns, including predators that dwell within the lake itself. The lake slopes sharply from its shores to a depth of 100 feet.

Four chuuls lurk near the lakebed, 90 feet below the surface. They're content to creep along the lightless depths and prey on cave fish unless something else catches their attention. Bright light draws the chuuls to the water's edge, as does any magic in the lake or near its surface, such as the boat from area L12. If drawn to the surface, the chuuls attack.

\subparagraph{L13a: Stone Bridge}
An unadorned stone bridge spans the northwest river.

\subparagraph{L13b: Gargoyle Bridge}
A bridge decorated with reliefs of gargoyles spans the west river flowing from the lake. The sound of roaring water echoes from downstream. A short distance beyond the bridge, the river becomes a waterfall, plummeting 300 feet into an Underdark lake. The lake and the creatures that inhabit it are beyond the scope of this adventure, but you can flesh them out as you see fit.

\subsubsection{L14: Basilisk Den}
\begin{DndReadAloud}{}
Oddly shaped chunks of rock are strewn about this cave, which reeks of ammonia. Scattered among the irregular rock pieces are stone sculptures of various creatures, including bats, subterranean lizards, and a broken bust of a gnome with a horrified expression.

\end{DndReadAloud}

Four basilisks dwell in this cave. Toward the back of the cave is a nest where the basilisks lay eggs. The basilisks are protective of their den, attacking any creatures that enter the cave.

\subparagraph{Treasure}
The head and shoulder of a petrified gnome lie near the basilisk nest at the back of the cave. The broken statue wears a pair of unpetrified goggles, which are \textit{Eyes of Minute Seeing}.

\subsubsection{L15: Rainbow Cavern}
\begin{DndReadAloud}{}
This vaulted cavern displays a rainbow of colors on its walls and floor. Stalactites hang like colorful icicles, and the stalagmites feature bright streaks and swirls. Mineral deposits form frozen curtains, cascades, and other fantastical shapes in vibrant hues.

\end{DndReadAloud}

The ceiling of this colorful cavern rises to a height of 50 feet. Careful examination of the floor reveals many cracked bone fragments among the stalagmites. Once the lair of a chimera, this cavern is now home to Lludd, a clever and long-lived \textbf{behir} who ousted the smaller, weaker creature that preceded him. Lludd speaks Common and Giant in addition to Draconic. The behir lurks on a rufous ledge near the entrance to the chamber, waiting for morsels worthy of his hunger.

Lludd is hostile toward intruders, but he doesn't immediately resort to violence. Though he's selfish and cruel, the behir hasn't survived this long by attempting to devour every creature that enters his lair—especially hardened adventurers. When he notices the characters, Lludd greets them in Common.

\subparagraph{Parleying with Lludd}
If the characters converse with Lludd, he offers to point them toward a place where they can find treasure in exchange for leaving him in peace. The behir keeps no hoard of his own. If they agree, Lludd directs them toward a "vast treasure" in the black-water lake (area L18).

During the conversation, Lludd sizes up the party. If the characters appear wounded or offend the behir, he seizes his chance and unleashes his breath on them.

\subparagraph{Treasure}
Inside Lludd's stomach is a \textit{Periapt of Proof against Poison} left over from an unlucky explorer the behir digested.

\subsubsection{L16: Boulder Heap}
\begin{DndReadAloud}{}
Small, rounded boulders are carefully stacked here, blocking the entrance to a wide tunnel.

\end{DndReadAloud}

The boulders are obviously not a natural formation. Lludd, the behir in area L15, piled them here to block the entrance to the greater caverns and prevent creatures that dwell there from taking him by surprise.

With 5 minutes of work, the characters can clear enough of the rocks to expose the tunnel beyond. There, a flight of nine hundred crudely shaped stone steps descends to the greater caverns. If the characters descend the steps, proceed to this adventure's conclusion.

\subparagraph{Hollow Boulder}
If the characters move any of the rocks, they uncover a smooth, blue-green boulder that's strangely light and clunks dully when moved or shaken, revealing its hollowness. The boulder can be easily broken open with blunt weapons, with tools, or by smashing it against a hard surface.

\subparagraph{Graven Glyphs}
Inside the hollow boulder is a bronze tablet worth 100 gp. The tablet is etched with a poem for those hoping to delve into the greater caverns. It offers cryptic insight on how to enter Iggwilv's inner sanctum in the greater caverns, where the archmage's true prize resides. The poem reads as follows:

\begin{DndReadAloud}{}
In the center lies the gate, But its locks will surely vex. Many are the guards who wait As you seek the middle hex.

Randomly sent to find a way Back to a different iron door. A seventh time and you may stay, And seek the glowing prize no more.

You have won old Iggwilv's prize, Her hoarded cache of magic, And freed the one with yearning eyes, Whose lot was hunger tragic.

\end{DndReadAloud}

\subsubsection{L17: Lavish Parlor}
\begin{DndReadAloud}{}
The natural cavern gives way to a luxurious parlor. Crystal chandeliers bathe the chamber in soft light, and the scent of orange blossoms perfumes the air. Plump cushions surround a low table heaped with ripe fruits, refreshing beverages, and succulent meats. Low sofas and comfortable chairs dot the walls of the lounge, and side tables between these sitting areas overflow with gemstones, jewelry, and other treasures.

Four impeccably dressed human aristocrats laugh and feast in the parlor. An unseen string quartet accompanies their revelry with soothing music.

\end{DndReadAloud}

This area is brightly lit. The nobles invite the characters to join the repast, insisting the characters rest and regale them with tales of adventure.

The nobles are part of an elaborate, genie-wrought illusion created by Kashem, a bitter \textbf{dao} bound to guard this chamber by Iggwilv. The Witch Queen imparted a mote of her power to Kashem in exchange for the dao's services. Long after the archmage's departure, Kashem still pursues her original charge to destroy all who enter.

When the characters arrive, Kashem floats silently above the chamber, waiting for a prime opportunity to strike. If the characters consume any of the food or drink, refuse to join the nobles outright, or see through the illusions, Kashem reveals herself and attacks.

\subparagraph{Illusory Finery}
The room's genie-wrought furnishings and inhabitants are convincing illusions. These illusions blanket dingy, neglected versions of the objects they cover and feel real if touched. To pierce the dao's ruse, a character examining the chamber's fixtures must succeed on a DC 20 Intelligence (Investigation) check. The nobles, however, fail to hold up to physical inspection. A character who touches a noble automatically succeeds on this check.

\subparagraph{Poisoned Feast}
The food and drink on the table are real—and poisoned. A creature that consumes any of the food or drink must make a DC 16 Constitution saving throw. On a failed save, a creature takes 21 (6d6) poison damage and has the poisoned condition for 1 hour. On a successful save, the creature takes half as much damage only. The illusory nobles nonchalantly partake of the sumptuous feast, politely encouraging the characters to sample the dishes.

\subparagraph{Dust to Dust}
If Kashem is slain, her body disintegrates into a pile of fine dust. This dust can be used in area L20 to destroy the cage in the lake or prove that Kashem is dead.

\subparagraph{Treasure}
None of the treasure is real. Coins and precious objects taken from the chamber turn out to be flat stones and obvious counterfeits.

\subsubsection{L18: Black-Water Lake}
\begin{DndReadAloud}{}
The river flows into a flooded cavern, creating a small lake. The surface of the water is glossy and black. A rocky island rises from the center of the pool.

\end{DndReadAloud}

The island is strewn with tiny garnets, which sparkle alluringly in the presence of light.

\subparagraph{Lurking Death}
Eight aquatic ghasts (each has a swimming speed of 30 feet) lurk in the lake, which is 20 feet deep. Spotting them beneath the water's mirrorlike surface requires a successful DC 20 Wisdom (Perception) check. The ghasts ambush any creatures that enter the lake.

\subparagraph{Rebuking Island}
Iggwilv recorded a series of messages in this cave to taunt intruders. Each time a creature sets foot on the island, an imperious voice echoes throughout the cave, uttering one of the following reproofs:


\begin{itemize}
\item "Fools! This is a dead end. Flee while you still can!"
\item "You were witless to enter at all, so you should probably just stay here."
\item (Peals of hideous, mocking laughter)
\end{itemize}

\subparagraph{Treasure}
Six hundred garnets are strewn across the island, each about the size of a pea and worth 1 gp. It takes 10 minutes for a single creature to collect the gems by hand. The gems have a combined weight of 5 pounds.

\subsubsection{L19: Cave of Crystals}
\begin{DndReadAloud}{}
Veins of glittering quartz streak the walls of this crystal cave. Three bizarre creatures with stone-lidded eyes converse in grating tones. Each has three clawed arms, three stumpy legs, and a barrel-like body.

\end{DndReadAloud}

Three \textbf{xorn} from the Elemental Plane of Earth found their way to this cave through a planar rift that has since closed, stranding them here. The xorn are initially indifferent to the characters but defend themselves if attacked.

\subparagraph{Hungry Castaways}
The xorn are discussing their plight when the characters arrive. The creatures are hungry and upset at their situation. Prone to stress eating, the xorn crave any precious gems or metals they can sense on the characters. If anyone in the party can understand the xorn, the creatures explain their dilemma and attempt to persuade the characters to give them a meal of at least 100 gp worth of gems or metals per xorn.

If the characters have a means to send the xorn home—such as the \textit{Banishment} spell—the xorn offer to help locate other treasure for the characters as payment. Alternatively, a single character can convince the xorn to assist them in defeating a threat elsewhere in the caverns with a successful DC 18 Charisma (Persuasion) check.

\subparagraph{Treasure}
The xorn stashed 500 gp worth of uncut gems behind a large crystal formation in the northwest corner of the room. A character searching the crystals finds the stash with a successful DC 17 Wisdom (Perception) check.

\subsubsection{L20: Pool Cavern}
\begin{DndReadAloud}{}
The narrow river flows into a domed cavern whose ceiling bristles with stalactites. Faint sounds of trickling water echo from the northwest.

\end{DndReadAloud}

The river dead-ends in this flooded chamber. A \textbf{marid} named Kasdu'ul is trapped at the bottom of the pool, 30 feet below the surface. An unbreakable stone cage etched with glowing runes confines the genie to the pool. Characters who look into the pool and succeed on a DC 15 Wisdom (Perception) check notice the cage's faint, glimmering light in the water.

\subparagraph{Bound Genie}
Kasdu'ul was trapped here by the dao Kashem (see area L17) centuries ago, shortly before Kashem was fooled and bound by Iggwilv in turn. A permanent \textit{Antipathy/Sympathy} spell (save DC 18) protects the cage and repels creatures that aren't native to the Elemental Plane of Earth with its Antipathy effect. The runes on the cage invoke the hatred of Ogrémoch, the Prince of Evil Earth.

While the cage is intact, Kasdu'ul has the incapacitated condition and speaks in a lethargic tone. If the characters overcome the Antipathy effect and approach the cage, the genie drowsily asks them to destroy it.

\subparagraph{Destroying the Cage}
The cage is immune to all damage and has no locks or opening mechanisms. It can be destroyed in the following ways:


\begin{itemize}
\item A creature native to the Plane of Earth—such as one of the pechs in area L5 or the xorn in area L19—can easily break the cage as if it were made of fragile clay.
\item If the dust of a slain dao is scattered over the lake, the cage instantly rusts away.
\item A successful \textit{Dispel Magic} (DC 18) spell cast on the cage causes it to crumble.
\end{itemize}

Kasdu'ul has had many years to contemplate the cage's magical construction. She might propose one of these methods to the party from inside the cage.

\subparagraph{Freed Genie}
If the cage is destroyed, Kasdu'ul returns to her lively self. The grateful genie urges the characters to find and slay the dao that imprisoned her if they haven't already.

As thanks for freeing her, she's willing to cast \textit{Tongues}, \textit{Water Breathing}, and \textit{Water Walk} on the characters to facilitate their exploration. If the characters also present proof of Kashem's demise—such as a pinch of the dust left behind after her death—or they offer a convincing story and succeed on a DC 17 Charisma (Deception) check, Kasdu'ul summons a \textbf{water weird} and orders it to serve the characters for 1 hour, after which it vanishes. In either case, Kasdu'ul then bids the characters farewell and returns to the Elemental Plane of Water.

\section{Conclusion}
Whether the characters found the entrance to the greater caverns (see area L16), succumbed to the dungeon's many threats, wisely decided to turn back, or simply ran out of time, Iggwilv's greatest treasures remain yet to be discovered. This adventure ends on a cliffhanger.

The rest of the Witch Queen's lair, her hoard, and even more iconic adventures can be found in Quests from the Infinite Staircase!

\appendix
\section{Magic Items}
This section presents four magic items related to Iggwilv and the Lost Caverns. These items are presented in alphabetical order.


\begin{itemize}
\begin{multicols}{3}
\item \textit{Demon Skin}
\item \textit{Iggwilv's Horn}
\item \textit{Potion of Polychromy}
\item \textit{Tasha's Creeping Keelboat}
\end{multicols}

\end{itemize}

\section{Creature}

\begin{itemize}
\item \textbf{Pech}
\end{itemize}

\section{Score Sheet}
Use this score sheet to tally the players' points. Players can gain or lose points for each of these actions only once.

\begin{DndTable}[header=Points by Area]{XXc}

 & Action & Points\\
L1: Entry Caverns & Caused a face to say "aah" & +1\\
 & Asked a face which way to go or used the word "truth" & +1\\
 & Recovered any gems & +1\\
 & Got bit by a stone face & −1\\
L2: Slate Chamber & Defeated the clay golem & +1\\
 & Removed the sword stuck in the clay golem's chest & +1\\
L3: Guano-Covered Cave & Defeated the cloaker & +1\\
 & Used the tunnel ledge to access the areas to the north & +1\\
L4: Littered Cave & Avoided the crushing ceiling trap & +1\\
 & Recovered the treasure & +1\\
L5: Grotto & Attacked the pechs & −1\\
 & Improved the pechs' mood & +2\\
 & Refused to douse any bright lights & −1\\
L6: Fungal Cavern & Disturbed the crickets & −1\\
L7: Slimy Cavern & Gathered or ate contaminated mushrooms & −1\\
 & Recovered the elf corpse's equipment & +1\\
L8: Stinking Cave & Defeated the trolls & +3\\
 & Found the trolls' loot & +1\\
L9: Bat Corridor & Disturbed the bats & −1\\
L10: Long Gallery & Attracted the attention of the fomorians in area L11 & −1\\
 & Pulled a cave moray from its burrow & +1\\
L11: Fomorian Cave & Defeated the fomorians & +3\\
 & Found the fomorians' loot & +1\\
L12: River Ledge & Learned the boat's command word & +1\\
L13: Underground Lake & Defeated the chuuls & +2\\
 & Failed a Strength saving throw in any of the rivers & −1\\
 & Cast \textit{Control Water}, \textit{Water Breathing}, or \textit{Water Walk} & +1\\
L14: Basilisk Den & Defeated the basilisks & +1\\
 & One or more characters gained the petrified condition & −1\\
 & Recovered the treasure near the basilisk's nest & +1\\
L15: Rainbow Cavern & Defeated Lludd the behir & +3\\
L16: Boulder Heap & Cleared the rocks to expose the tunnel to the greater caverns & +1\\
 & Discovered the bronze tablet within the hollow boulder & +1\\
L17: Lavish Parlor & Ate any of the poisoned food or drink & −1\\
 & Defeated Kashem the dao & +2\\
 & Saw through Kashem's illusions before she attacked & +1\\
L18: Black-Water Lake & Gathered all the garnets on the island & +1\\
 & Received an insult from Iggwilv & +1\\
 & Used a cleric or paladin's Channel Divinity feature against the ghasts & +1\\
L19: Cave of Crystals & Attacked the xorn & −1\\
 & Sent any xorn home & +2\\
 & Enlisted the help of the xorn & +1\\
L20: Pool Caverns & Freed Kasdu'ul the marid & +2\\
 & Provided proof of Kashem's demise or convinced Kasdu'ul of the dao's defeat & +1\\
\end{DndTable}

\begin{DndTable}[header=Additional Points]{Xc}

Condition & Points\\
Per character death & −1\\
One or more players thanked the DM during or after the session & +1\\
\end{DndTable}

\section{Premade Characters}
If a player doesn't have a level-appropriate character to play, have them choose one from this section. The Premade Characters table summarizes the characters, all of whom are 9th level. Give players who choose from these options time to read over their character sheets before starting the adventure.

\begin{DndTable}[header=Premade Characters]{XXX}

Name & Species & Class\\
Cathartic & Human & Cleric\\
Dunil & Halfling & Rogue\\
Ethelrede & Human & Fighter\\
Flemin & Dwarf & Monk\\
Rustle Berrydust & Gnome & Wizard\\
Weslocke & Tiefling & Warlock\\
\end{DndTable}

\section{Credits}

\begin{description}
\begin{multicols}{2}

\begin{description}
\item[Lead Designer.]
Justice Ramin Arman

\item[Designer.]
Dan Dillon

\item[Rules Developer.]
Ron Lundeen

\item[Art Director.]
Fury Galluzzi

\item[Lead Editor.]
Judy Bauer

\item[Editor.]
Hannah Rose

\item[Graphic Designer.]
Paolo Vacala

\item[Cover Illustrator.]
Jodie Muir

\item[Cartographer.]
Mike Schley

\item[Interior Illustrators.]
Stephen Andrade, Mark Behm, Zoltan Boros, Adrián Ibarra Lugo, Arash Radkia

\item[Premade Character Portraits.]
Fay Dalton

\item[Consultants.]
James Mendez-Hodes, Pam Punzalan

\item[Game Architects.]
Jeremy Crawford, Christopher Perkins

\item[Studio Art Director.]
Josh Herman

\item[Senior Producer.]
Dan Tovar

\item[Producers.]
Bill Benham, Siera Bruggeman, Robert Hawkey

\item[Product Manager.]
Natalie Egan

\end{description}

\section{D\&D Beyond}

\begin{description}
\item[Product Manager.]
Jeff Turriff

\item[Digital Design Team.]
Jay Jani, Sean Stoves, Adam Walton

\end{description}

\end{multicols}

\end{description}

\subsection{Original Credits}

\begin{description}
\item[Design.]
Gary Gygax

\item[Development.]
Gary Gygax, Allen Hammack, Jon Pickens, Edward G. Sollers

\item[Editing.]
Edward G. Sollers

\item[Art.]
Jim Holloway, Erol Otus, Jeff Easley, Stephen D. Sullivan

\item[Playtesting.]
Jeff Dolphin, Luke Gygax, David Kuntz, Richard Kuntz, Sonny Savage, James M. Ward

\item[Special Thanks.]
Rob Kuntz

\end{description}


\begin{description}
\item[Special thanks to the original designer of this adventure, Gary Gygax, as well as everyone who contributed to it:]
Jeff Dolphin, Luke Gygax, Jeff Easley, Allen Hammack, Jim Holloway, David Kuntz, Richard Kuntz, Erol Otus, Jon Pickens, Sonny Savage, Edward G. Sollers, Stephen D. Sullivan, and James M. Ward.
\end{description}

\chapter{Magic Items}
\DndItemHeader{Bracers of Defense}{r, a, r, e None}
While wearing these bracers, you gain a +2 bonus to AC if you are wearing no armor and using no shield.

\DndItemHeader{Eyes of Minute Seeing}{u, n, c, o, m, m, o, n}
These crystal lenses fit over the eyes. While wearing them, you can see much better than normal out to a range of 1 foot. You have advantage on Intelligence (Investigation) checks that rely on sight while searching an area or studying an object within that range.

\DndItemHeader{Folding Boat}{r, a, r, e}
This object appears as a wooden box that measures 12 inches long, 6 inches wide, and 6 inches deep. It weighs 4 pounds and floats. It can be opened to store items inside. This item also has three command words, each requiring you to use an action to speak it.

One command word causes the box to unfold into a boat 10 feet long, 4 feet wide, and 2 feet deep. The boat has one pair of oars, an anchor, a mast, and a lateen sail. The boat can hold up to four Medium creatures comfortably.

The second command word causes the box to unfold into a ship 24 feet long, 8 feet wide; and 6 feet deep. The ship has a deck, rowing seats, five sets of oars, a steering oar, an anchor, a deck cabin, and a mast with a square sail. The ship can hold fifteen Medium creatures comfortably.

When the box becomes a vessel, its weight becomes that of a normal vessel its size, and anything that was stored in the box remains in the boat.

The third command word causes the folding boat to fold back into a box, provided that no creatures are aboard. Any objects in the vessel that can't fit inside the box remain outside the box as it folds. Any objects in the vessel that can fit inside the box do so.

\DndItemHeader{Iggwilv's Horn}{r, a, r, e None}
You can use an action to blow this horn to cast one of the following spells from it: \textit{Arms of Hadar}, \textit{Fog Cloud}, \textit{Gust of Wind}, or \textit{Web}. If the spell requires a saving throw, the spell save DC is 13.

Once the horn has been used to cast a spell, it can't be used to cast that spell again until the next dawn.

\DndItemHeader{Periapt of Proof against Poison}{r, a, r, e}
This delicate silver chain has a brilliant-cut black gem pendant. While you wear it, poisons have no effect on you. You are immune to the poisoned condition and have immunity to poison damage.

\DndItemHeader{Potion of Greater Healing}{u, n, c, o, m, m, o, n}
You regain 4d4 + 4 hit points when you drink this potion. The potion's red liquid glimmers when agitated.

\DndItemHeader{Potion of Polychromy}{u, n, c, o, m, m, o, n}
When you drink this potion, you and everything you are wearing or carrying take on a rainbow-hued appearance for 1 hour. During that time, you can use a bonus action to turn any color or combination of colors you choose. If you mimic the colors of your surroundings, your hues continually shift to match your surroundings, and you have advantage on Dexterity (Stealth) checks until you change your colors again or the potion wears off.

The potion is separated into seven brightly colored bands of immiscible liquids and has a syrupy taste.

\DndItemHeader{Tasha's Creeping Keelboat}{v, e, r, y,  , r, a, r, e None}
This magic vehicle is a boat 10 feet wide and 30 feet long. It has four legs that propel it across land and water. It has a walking and swimming speed of 20 feet, but it can't travel underwater. The boat moves according to your spoken directions while you are riding it, and creatures of your choice gain a +1 bonus to their Armor Class while on the boat.

The boat can transport up to 1,000 pounds without hindrance. It can carry up to twice this weight, but it moves at half speed if it carries more than its normal capacity.

\DndItemHeader{Quiver}{n, o, n, e}
A quiver can hold up to 20 arrow.

\chapter{Creatures}
\setlist[description]{nolistsep,listparindent=3pt,topsep=6pt,font={\bfseries\sffamily\small}}
\pagestyle{empty}

\begin{DndMonster}[float=t]{Basilisk}
\DndMonsterType{Medium monstrosity, unaligned}
\DndMonsterBasics[
armor-class = {15 (natural armor)},
hit-points = {\DndDice{8d8 + 16}},
speed = {20 ft.}
]
\DndMonsterAbilityScores[
 str = 16,
 dex = 8,
 con = 15,
 int = 2,
 wis = 8,
 cha = 7
]
\DndMonsterDetails[
senses = {darkvision 60 ft., passive perception 9},
challenge = 3,
]
\DndMonsterAction{Petrifying Gaze}
If a creature starts its turn within 30 feet of the basilisk and the two of them can see each other, the basilisk can force the creature to make a DC 12 Constitution saving throw if the basilisk isn't incapacitated. On a failed save, the creature magically begins to turn to stone and is restrained. It must repeat the saving throw at the end of its next turn. On a success, the effect ends. On a failure, the creature is petrified until freed by the  \textit{greater restoration} spell or other magic.

A creature that isn't Surprise can avert its eyes to avoid the saving throw at the start of its turn. If it does so, it can't see the basilisk until the start of its next turn, when it can avert its eyes again. If it looks at the basilisk in the meantime, it must immediately make the save.

If the basilisk sees its reflection within 30 feet of it in bright light, it mistakes itself for a rival and targets itself with its gaze.

\DndMonsterSection{Actions}
\DndMonsterAction{Bite}
\textsl{Melee Weapon Attack:} 5 to hit, reach 5 ft., one target. \textsl{Hit:} 10 (2d6 + 3) piercing damage plus 7 (2d6) poison damage.

\end{DndMonster}



\begin{DndMonster}[float=t]{Behir}
\DndMonsterType{Huge monstrosity, neutral evil}
\DndMonsterBasics[
armor-class = {17 (natural armor)},
hit-points = {\DndDice{16d12 + 64}},
speed = {50 ft., climb 40 ft.}
]
\DndMonsterAbilityScores[
 str = 23,
 dex = 16,
 con = 18,
 int = 7,
 wis = 14,
 cha = 12
]
\DndMonsterDetails[
skills = {Perception +6, Stealth +7},
damage-immunities = {lightning},
senses = {darkvision 90 ft., passive perception 16},
languages = {Draconic},
challenge = 11,
]
\DndMonsterSection{Actions}
\DndMonsterAction{Multiattack}
The behir makes two attacks: one with its bite and one to constrict.

\DndMonsterAction{Bite}
\textsl{Melee Weapon Attack:} 10 to hit, reach 10 ft., one target. \textsl{Hit:} 22 (3d10 + 6) piercing damage.

\DndMonsterAction{Constrict}
\textsl{Melee Weapon Attack:} 10 to hit, reach 5 ft., one Large or smaller creature. \textsl{Hit:} 17 (2d10 + 6) bludgeoning damage plus 17 (2d10 + 6) slashing damage. The target is grappled (escape DC 16) if the behir isn't already constricting a creature, and the target is restrained until this grapple ends.

\DndMonsterAction{Lightning Breath (Recharge 5\textendash 6)}
The behir exhales a line of lightning that is 20 feet long and 5 feet wide. Each creature in that line must make a DC 16 Dexterity saving throw, taking 66 (12d10) lightning damage on a failed save, or half as much damage on a successful one.

\DndMonsterAction{Swallow}
The behir makes one bite attack against a Medium or smaller target it is grappling. If the attack hits, the target is also swallowed, and the grapple ends. While swallowed, the target is blinded and restrained, it has total cover against attacks and other effects outside the behir, and it takes 21 (6d6) acid damage at the start of each of the behir's turns. A behir can have only one creature swallowed at a time.

If the behir takes 30 damage or more on a single turn from the swallowed creature, the behir must succeed on a DC 14 Constitution saving throw at the end of that turn or regurgitate the creature, which falls prone in a space within 10 feet of the behir. If the behir dies, a swallowed creature is no longer restrained by it and can escape from the corpse by using 15 feet of movement, exiting prone.

\end{DndMonster}



\begin{DndMonster}[float=t]{Chuul}
\DndMonsterType{Large aberration, chaotic evil}
\DndMonsterBasics[
armor-class = {16 (natural armor)},
hit-points = {\DndDice{11d10 + 33}},
speed = {30 ft., swim 30 ft.}
]
\DndMonsterAbilityScores[
 str = 19,
 dex = 10,
 con = 16,
 int = 5,
 wis = 11,
 cha = 5
]
\DndMonsterDetails[
skills = {Perception +4},
damage-immunities = {poison},
condition-immunities = {poisoned},
senses = {darkvision 60 ft., passive perception 14},
languages = {understands Deep Speech but can't speak},
challenge = 4,
]
\DndMonsterAction{Amphibious}
The chuul can breathe air and water.

\DndMonsterAction{Sense Magic}
The chuul senses magic within 120 feet of it at will. This trait otherwise works like the \textit{detect magic} spell but isn't itself magical.

\DndMonsterSection{Actions}
\DndMonsterAction{Multiattack}
The chuul makes two pincer attacks. If the chuul is grappling a creature, the chuul can also use its tentacles once.

\DndMonsterAction{Pincer}
\textsl{Melee Weapon Attack:} 6 to hit, reach 10 ft., one target. \textsl{Hit:} 11 (2d6 + 4) bludgeoning damage. The target is grappled (escape DC 14) if it is a Large or smaller creature and the chuul doesn't have two other creatures grappled.

\DndMonsterAction{Tentacles}
One creature grappled by the chuul must succeed on a DC 13 Constitution saving throw or be poisoned for 1 minute. Until this poison ends, the target is paralyzed. The target can repeat the saving throw at the end of each of its turns, ending the effect on itself on a success.

\end{DndMonster}



\begin{DndMonster}[float=t]{Clay Golem}
\DndMonsterType{Large construct, unaligned}
\DndMonsterBasics[
armor-class = {14 (natural armor)},
hit-points = {\DndDice{14d10 + 56}},
speed = {20 ft.}
]
\DndMonsterAbilityScores[
 str = 20,
 dex = 9,
 con = 18,
 int = 3,
 wis = 8,
 cha = 1
]
\DndMonsterDetails[
damage-immunities = {acid; poison; psychic; bludgeoning, piercing, slashing from nonmagical attacks that aren't adamantine},
condition-immunities = {charmed, exhaustion, frightened, paralyzed, petrified, poisoned},
senses = {darkvision 60 ft., passive perception 9},
languages = {understands the languages of its creator but can't speak},
challenge = 9,
]
\DndMonsterAction{Acid Absorption}
Whenever the golem is subjected to acid damage, it takes no damage and instead regains a number of hit points equal to the acid damage dealt.

\DndMonsterAction{Berserk}
Whenever the golem starts its turn with 60 hit points or fewer, roll a d6. On a 6, the golem goes berserk. On each of its turns while berserk, the golem attacks the nearest creature it can see. If no creature is near enough to move to and attack, the golem attacks an object, with preference for an object smaller than itself. Once the golem goes berserk, it continues to do so until it is destroyed or regains all its hit points.

\DndMonsterAction{Immutable Form}
The golem is immune to any spell or effect that would alter its form.

\DndMonsterAction{Magic Resistance}
The golem has advantage on saving throws against spells and other magical effects.

\DndMonsterAction{Magic Weapons}
The golem's weapon attacks are magical.

\DndMonsterSection{Actions}
\DndMonsterAction{Multiattack}
The golem makes two slam attacks.

\DndMonsterAction{Slam}
\textsl{Melee Weapon Attack:} 8 to hit, reach 5 ft., one target. \textsl{Hit:} 16 (2d10 + 5) bludgeoning damage. If the target is a creature, it must succeed on a DC 15 Constitution saving throw or have its hit point maximum reduced by an amount equal to the damage taken. The target dies if this attack reduces its hit point maximum to 0. The reduction lasts until removed by the  \textit{greater restoration} spell or other magic.

\DndMonsterAction{Haste (Recharge 5\textendash 6)}
Until the end of its next turn, the golem magically gains a +2 bonus to its AC, has advantage on Dexterity saving throws, and can use its slam attack as a bonus action.

\end{DndMonster}



\begin{DndMonster}[float=t]{Cloaker}
\DndMonsterType{Large aberration, chaotic neutral}
\DndMonsterBasics[
armor-class = {14 (natural armor)},
hit-points = {\DndDice{12d10 + 12}},
speed = {10 ft., fly 40 ft.}
]
\DndMonsterAbilityScores[
 str = 17,
 dex = 15,
 con = 12,
 int = 13,
 wis = 12,
 cha = 14
]
\DndMonsterDetails[
skills = {Stealth +5},
senses = {darkvision 60 ft., passive perception 11},
languages = {Deep Speech, Undercommon},
challenge = 8,
]
\DndMonsterAction{Damage Transfer}
While attached to a creature, the cloaker takes only half the damage dealt to it (rounded down). and that creature takes the other half.

\DndMonsterAction{False Appearance}
While the cloaker remains motionless without its underside exposed, it is indistinguishable from a dark leather cloak.

\DndMonsterAction{Light Sensitivity}
While in bright light, the cloaker has disadvantage on attack rolls and Wisdom (Perception) checks that rely on sight.

\DndMonsterSection{Actions}
\DndMonsterAction{Multiattack}
The cloaker makes two attacks: one with its bite and one with its tail.

\DndMonsterAction{Bite}
\textsl{Melee Weapon Attack:} 6 to hit, reach 5 ft., one creature. \textsl{Hit:} 10 (2d6 + 3) piercing damage, and if the target is Large or smaller, the cloaker attaches to it. If the cloaker has advantage against the target, the cloaker attaches to the target's head, and the target is blinded and unable to breathe while the cloaker is attached. While attached, the cloaker can make this attack only against the target and has advantage on the attack roll. The cloaker can detach itself by spending 5 feet of its movement. A creature, including the target, can take its action to detach the cloaker by succeeding on a DC 16 Strength check.

\DndMonsterAction{Tail}
\textsl{Melee Weapon Attack:} 6 to hit, reach 10 ft., one creature. \textsl{Hit:} 7 (1d8 + 3) slashing damage.

\DndMonsterAction{Moan}
Each creature within 60 feet of the cloaker that can hear its moan and that isn't an aberration must succeed on a DC 13 Wisdom saving throw or become frightened until the end of the cloaker's next turn. If a creature's saving throw is successful, the creature is immune to the cloaker's moan for the next 24 hours.

\DndMonsterAction{Phantasms (Recharges after a Short or Long Rest)}
The cloaker magically creates three illusory duplicates of itself if it isn't in bright light. The duplicates move with it and mimic its actions, shifting position so as to make it impossible to track which cloaker is the real one. If the cloaker is ever in an area of bright light, the duplicates disappear.

Whenever any creature targets the cloaker with an attack or a harmful spell while a duplicate remains, that creature rolls randomly to determine whether it targets the cloaker or one of the duplicates. A creature is unaffected by this magical effect if it can't see or if it relies on senses other than sight.

A duplicate has the cloaker's AC and uses its saving throws. If an attack hits a duplicate, or if a duplicate fails a saving throw against an effect that deals damage, the duplicate disappears.

\end{DndMonster}



\begin{DndMonster}[float=t]{Dao}
\DndMonsterType{Large elemental, neutral evil}
\DndMonsterBasics[
armor-class = {18 (natural armor)},
hit-points = {\DndDice{15d10 + 105}},
speed = {30 ft., burrow 30 ft., fly 30 ft.}
]
\DndMonsterAbilityScores[
 str = 23,
 dex = 12,
 con = 24,
 int = 12,
 wis = 13,
 cha = 14
]
\DndMonsterDetails[
saving-throws = {Int +5, Wis +5, Cha +6},
condition-immunities = {petrified},
senses = {darkvision 120 ft., passive perception 11},
languages = {Terran},
challenge = 11,
]
\DndMonsterAction{Earth Glide}
The dao can burrow through nonmagical, unworked earth and stone. While doing so, the dao doesn't disturb the material it moves through.

\DndMonsterAction{Elemental Demise}
If the dao dies, its body disintegrates into crystalline powder, leaving behind only equipment the dao was wearing or carrying.

\DndMonsterAction{Sure-Footed}
The dao has advantage on Strength and Dexterity saving throws made against effects that would knock it prone.

\DndMonsterAction{Innate Spellcasting}
The dao's innate spellcasting ability is Charisma (spell save DC 14, 6 to hit with spell attacks). It can innately cast the following spells, requiring no material components:
\begin{DndMonsterSpells}
  \DndInnateSpellLevel{\textit{detect evil and good}, \textit{detect magic}, \textit{stone shape}}
  \DndInnateSpellLevel[3e]{\textit{passwall}, \textit{move earth}, \textit{tongues}}
  \DndInnateSpellLevel[1e]{\textit{conjure elemental} (\textbf{earth elemental} only), \textit{gaseous form}, \textit{invisibility}, \textit{phantasmal killer}, \textit{plane shift}, \textit{wall of stone}}
\end{DndMonsterSpells}
\DndMonsterSection{Actions}
\DndMonsterAction{Multiattack}
The Dao makes two fist attacks or two maul attacks.

\DndMonsterAction{Fist}
\textsl{Melee Weapon Attack:} 10 to hit, reach 5 ft., one target. \textsl{Hit:} 15 (2d8 + 6) bludgeoning damage.

\DndMonsterAction{Maul}
\textsl{Melee Weapon Attack:} 10 to hit, reach 5 ft., one target. \textsl{Hit:} 20 (4d6 + 6) bludgeoning damage. If the target is a Huge or smaller creature, it must succeed on a DC 18 Strength check or be knocked prone.

\end{DndMonster}



\begin{DndMonster}[float=t]{Earth Elemental}
\DndMonsterType{Large elemental, neutral}
\DndMonsterBasics[
armor-class = {17 (natural armor)},
hit-points = {\DndDice{12d10 + 60}},
speed = {30 ft., burrow 30 ft.}
]
\DndMonsterAbilityScores[
 str = 20,
 dex = 8,
 con = 20,
 int = 5,
 wis = 10,
 cha = 5
]
\DndMonsterDetails[
damage-vulnerabilities = {thunder},
damage-resistances = {bludgeoning, piercing, slashing from nonmagical attacks},
damage-immunities = {poison},
condition-immunities = {exhaustion, paralyzed, petrified, poisoned, unconscious},
senses = {darkvision 60 ft., tremorsense 60 ft., passive perception 10},
languages = {Terran},
challenge = 5,
]
\DndMonsterAction{Earth Glide}
The elemental can burrow through nonmagical, unworked earth and stone. While doing so, the elemental doesn't disturb the material it moves through.

\DndMonsterAction{Siege Monster}
The elemental deals double damage to objects and structures.

\DndMonsterSection{Actions}
\DndMonsterAction{Multiattack}
The elemental makes two slam attacks.

\DndMonsterAction{Slam}
\textsl{Melee Weapon Attack:} 8 to hit, reach 10 ft., one target. \textsl{Hit:} 14 (2d8 + 5) bludgeoning damage.

\end{DndMonster}



\begin{DndMonster}[float=t]{Fomorian}
\DndMonsterType{Huge giant, chaotic evil}
\DndMonsterBasics[
armor-class = {14 (natural armor)},
hit-points = {\DndDice{13d12 + 65}},
speed = {30 ft.}
]
\DndMonsterAbilityScores[
 str = 23,
 dex = 10,
 con = 20,
 int = 9,
 wis = 14,
 cha = 6
]
\DndMonsterDetails[
skills = {Perception +8, Stealth +3},
senses = {darkvision 120 ft., passive perception 18},
languages = {Giant, Undercommon},
challenge = 8,
]
\DndMonsterSection{Actions}
\DndMonsterAction{Multiattack}
The fomorian attacks twice with its greatclub or makes one greatclub attack and uses Evil Eye once.

\DndMonsterAction{Greatclub}
\textsl{Melee Weapon Attack:} 9 to hit, reach 15 ft., one target. \textsl{Hit:} 19 (3d8 + 6) bludgeoning damage.

\DndMonsterAction{Evil Eye}
The fomorian magically forces a creature it can see within 60 feet of it to make a DC 14 Charisma saving throw. The creature takes 27 (6d8) psychic damage on a failed save, or half as much damage on a successful one.

\DndMonsterAction{Curse of the Evil Eye (Recharges after a Short or Long Rest)}
With a stare, the fomorian uses Evil Eye, but on a failed save, the creature is also cursed with magical deformities. While deformed, the creature has its speed halved and has disadvantage on ability checks, saving throws, and attacks based on Strength or Dexterity.

The transformed creature can repeat the saving throw whenever it finishes a long rest, ending the effect on a success.

\end{DndMonster}



\begin{DndMonster}[float=t]{Ghast}
\DndMonsterType{Medium undead, chaotic evil}
\DndMonsterBasics[
armor-class = {13},
hit-points = {\DndDice{8d8}},
speed = {30 ft.}
]
\DndMonsterAbilityScores[
 str = 16,
 dex = 17,
 con = 10,
 int = 11,
 wis = 10,
 cha = 8
]
\DndMonsterDetails[
damage-resistances = {necrotic},
damage-immunities = {poison},
condition-immunities = {charmed, exhaustion, poisoned},
senses = {darkvision 60 ft., passive perception 10},
languages = {Common},
challenge = 2,
]
\DndMonsterAction{Stench}
Any creature that starts its turn within 5 feet of the ghast must succeed on a DC 10 Constitution saving throw or be poisoned until the start of its next turn. On a successful saving throw, the creature is immune to the ghast's Stench for 24 hours.

\DndMonsterAction{Turn Defiance}
The ghast and any ghouls within 30 feet of it have advantage on saving throws against effects that turn undead.

\DndMonsterSection{Actions}
\DndMonsterAction{Bite}
\textsl{Melee Weapon Attack:} 3 to hit, reach 5 ft., one creature. \textsl{Hit:} 12 (2d8 + 3) piercing damage.

\DndMonsterAction{Claws}
\textsl{Melee Weapon Attack:} 5 to hit, reach 5 ft., one target. \textsl{Hit:} 10 (2d6 + 3) slashing damage. If the target is a creature other than an undead, it must succeed on a DC 10 Constitution saving throw or be paralyzed for 1 minute. The target can repeat the saving throw at the end of each of its turns, ending the effect on itself on a success.

\end{DndMonster}



\begin{DndMonster}[float=t]{Giant Frog}
\DndMonsterType{Medium beast, unaligned}
\DndMonsterBasics[
armor-class = {11},
hit-points = {\DndDice{4d8}},
speed = {30 ft., swim 30 ft.}
]
\DndMonsterAbilityScores[
 str = 12,
 dex = 13,
 con = 11,
 int = 2,
 wis = 10,
 cha = 3
]
\DndMonsterDetails[
skills = {Perception +2, Stealth +3},
senses = {darkvision 30 ft., passive perception 12},
challenge = 1/4,
]
\DndMonsterAction{Amphibious}
The frog can breathe air and water.

\DndMonsterAction{Standing Leap}
The frog's long jump is up to 20 feet and its high jump is up to 10 feet, with or without a running start.

\DndMonsterSection{Actions}
\DndMonsterAction{Bite}
\textsl{Melee Weapon Attack:} 3 to hit, reach 5 ft., one target. \textsl{Hit:} 4 (1d6 + 1) piercing damage, and the target is grappled (escape DC 11). Until this grapple ends, the target is restrained, and the frog can't bite another target.

\DndMonsterAction{Swallow}
The frog makes one bite attack against a Small or smaller target it is grappling. If the attack hits, the target is swallowed, and the grapple ends. The swallowed target is blinded and restrained, it has total cover against attacks and other effects outside the frog, and it takes 5 (2d4) acid damage at the start of each of the frog's turns. The frog can have only one target swallowed at a time. If the frog dies, a swallowed creature is no longer restrained by it and can escape from the corpse using 5 feet of movement, exiting prone.

\end{DndMonster}



\begin{DndMonster}[float=t]{Giant Poisonous Snake}
\DndMonsterType{Medium beast, unaligned}
\DndMonsterBasics[
armor-class = {14},
hit-points = {\DndDice{2d8 + 2}},
speed = {30 ft., swim 30 ft.}
]
\DndMonsterAbilityScores[
 str = 10,
 dex = 18,
 con = 13,
 int = 2,
 wis = 10,
 cha = 3
]
\DndMonsterDetails[
skills = {Perception +2},
senses = {blindsight 10 ft., passive perception 12},
challenge = 1/4,
]
\DndMonsterSection{Actions}
\DndMonsterAction{Bite}
\textsl{Melee Weapon Attack:} 6 to hit, reach 10 ft., one target. \textsl{Hit:} 6 (1d4 + 4) piercing damage, and the target must make a DC 11 Constitution saving throw, taking 10 (3d6) poison damage on a failed save, or half as much damage on a successful one.

\end{DndMonster}



\begin{DndMonster}[float=t]{Marid}
\DndMonsterType{Large elemental, chaotic neutral}
\DndMonsterBasics[
armor-class = {17 (natural armor)},
hit-points = {\DndDice{17d10 + 136}},
speed = {30 ft., fly 60 ft., swim 90 ft.}
]
\DndMonsterAbilityScores[
 str = 22,
 dex = 12,
 con = 26,
 int = 18,
 wis = 17,
 cha = 18
]
\DndMonsterDetails[
saving-throws = {Dex +5, Wis +7, Cha +8},
damage-resistances = {acid, cold, lightning},
senses = {blindsight 30 ft., darkvision 120 ft., passive perception 13},
languages = {Aquan},
challenge = 11,
]
\DndMonsterAction{Amphibious}
The marid can breathe air and water.

\DndMonsterAction{Elemental Demise}
If the marid dies, its body disintegrates into a burst of water and foam, leaving behind only equipment the marid was wearing or carrying.

\DndMonsterAction{Innate Spellcasting}
The marid's innate spellcasting ability is Charisma (spell save DC 16, 8 to hit with spell attacks). It can innately cast the following spells, requiring no material components:
\begin{DndMonsterSpells}
  \DndInnateSpellLevel{\textit{create or destroy water}, \textit{detect evil and good}, \textit{detect magic}, \textit{fog cloud}, \textit{purify food and drink}}
  \DndInnateSpellLevel[3e]{\textit{tongues}, \textit{water breathing}, \textit{water walk}}
  \DndInnateSpellLevel[1e]{\textit{conjure elemental} (\textbf{water elemental} only), \textit{control water}, \textit{gaseous form}, \textit{invisibility}, \textit{plane shift}}
\end{DndMonsterSpells}
\DndMonsterSection{Actions}
\DndMonsterAction{Multiattack}
The marid makes two trident attacks.

\DndMonsterAction{Trident}
\textsl{Melee or Ranged Weapon Attack:} 10 to hit, reach 5 ft. or range 20/60 ft., one target. \textsl{Hit:} 13 (2d6 + 6) piercing damage, or 15 (2d8 + 6) piercing damage if used with two hands to make a melee attack.

\DndMonsterAction{Water Jet}
The marid magically shoots water in a 60-foot line that is 5 feet wide. Each creature in that line must make a DC 16 Dexterity saving throw. On a failure, a target takes 21 (6d6) bludgeoning damage and, if it is Huge or smaller, is pushed up to 20 feet away from the marid and knocked prone. On a success, a target takes half the bludgeoning damage, but is neither pushed nor knocked prone.

\end{DndMonster}



\begin{DndMonster}[float=t]{Pech}
\DndMonsterType{Small elemental, neutral good}
\DndMonsterBasics[
armor-class = {17 (natural armor)},
hit-points = {\DndDice{11d6 + 44}},
speed = {30 ft., burrow 20 ft.}
]
\DndMonsterAbilityScores[
 str = 19,
 dex = 11,
 con = 18,
 int = 11,
 wis = 14,
 cha = 10
]
\DndMonsterDetails[
saving-throws = {Con +6, Wis +4},
skills = {Athletics +6, Perception +4, Survival +4},
condition-immunities = {petrified},
senses = {darkvision 120 ft., tremorsense 120 ft., passive perception 14},
languages = {Common, Terran},
challenge = 4,
]
\DndMonsterAction{Magic Resistance}
The pech has advantage on saving throws against spells and other magical effects.

\DndMonsterAction{Sunlight Sensitivity}
While in sunlight, the pech has disadvantage on attack rolls.

\DndMonsterAction{Communal Spellcasting (2/Day)}
The pech works with three or more pechs to cast spells, requiring no spell components and using Wisdom as the spellcasting ability (save DC 12). If at least three other pechs are within 30 feet of it, the pech can cast \textit{Wall of Stone}. If at least seven other pechs are within 30 feet of it, it can cast \textit{Greater Restoration}. Each other pech involved in casting the spell can't have the incapacitated condition and must have at least one use of Communal Spellcasting remaining, which it must immediately expend to participate (no action required).
\begin{DndMonsterSpells}
  \DndInnateSpellLevel[2]{\textit{wall of stone}, \textit{greater restoration}}
\end{DndMonsterSpells}
\DndMonsterAction{Stone Shape (3/Day)}
The pech casts \textit{Stone Shape}, requiring no spell components and using Wisdom as the spellcasting ability.
\begin{DndMonsterSpells}
  \DndInnateSpellLevel[3e]{\textit{stone shape}}
\end{DndMonsterSpells}
\DndMonsterSection{Actions}
\DndMonsterAction{Multiattack}
The pech makes two Fortified Pickaxe attacks. If it hits a Large or smaller creature with both attacks, the target must succeed on a DC 14 Strength saving throw or have the prone condition.

\DndMonsterAction{Fortified Pickaxe}
\textsl{Melee Weapon Attack:} 6 to hit, reach 5 ft., one target. \textsl{Hit:} 11 (2d6 + 4) force damage. If the target is a Construct or an object, the attack is automatically a critical hit.

\end{DndMonster}



\begin{DndMonster}[float=t]{Swarm of Bats}
\DndMonsterType{Medium beast, unaligned}
\DndMonsterBasics[
armor-class = {12},
hit-points = {\DndDice{5d8}},
speed = {0 ft., fly 30 ft.}
]
\DndMonsterAbilityScores[
 str = 5,
 dex = 15,
 con = 10,
 int = 2,
 wis = 12,
 cha = 4
]
\DndMonsterDetails[
damage-resistances = {bludgeoning, piercing, slashing},
condition-immunities = {charmed, frightened, grappled, paralyzed, petrified, prone, restrained, stunned},
senses = {blindsight 60 ft., passive perception 11},
challenge = 1/4,
]
\DndMonsterAction{Echolocation}
The swarm can't use its blindsight while deafened.

\DndMonsterAction{Keen Hearing}
The swarm has advantage on Wisdom (Perception) checks that rely on hearing.

\DndMonsterAction{Swarm}
The swarm can occupy another creature's space and vice versa, and the swarm can move through any opening large enough for a Tiny bat. The swarm can't regain hit points or gain temporary hit points.

\DndMonsterSection{Actions}
\DndMonsterAction{Bites}
\textsl{Melee Weapon Attack:} 4 to hit, reach 0 ft., one creature in the swarm's space. \textsl{Hit:} 5 (2d4) piercing damage, or 2 (1d4) piercing damage if the swarm has half of its hit points or fewer.

\end{DndMonster}



\begin{DndMonster}[float=t]{Troll}
\DndMonsterType{Large giant, chaotic evil}
\DndMonsterBasics[
armor-class = {15 (natural armor)},
hit-points = {\DndDice{8d10 + 40}},
speed = {30 ft.}
]
\DndMonsterAbilityScores[
 str = 18,
 dex = 13,
 con = 20,
 int = 7,
 wis = 9,
 cha = 7
]
\DndMonsterDetails[
skills = {Perception +2},
senses = {darkvision 60 ft., passive perception 12},
languages = {Giant},
challenge = 5,
]
\DndMonsterAction{Keen Smell}
The troll has advantage on Wisdom (Perception) checks that rely on smell.

\DndMonsterAction{Regeneration}
The troll regains 10 hit points at the start of its turn. If the troll takes acid or fire damage, this trait doesn't function at the start of the troll's next turn. The troll dies only if it starts its turn with 0 hit points and doesn't regenerate.

\DndMonsterSection{Actions}
\DndMonsterAction{Multiattack}
The troll makes three attacks: one with its bite and two with its claws.

\DndMonsterAction{Bite}
\textsl{Melee Weapon Attack:} 7 to hit, reach 5 ft., one target. \textsl{Hit:} 7 (1d6 + 4) piercing damage.

\DndMonsterAction{Claw}
\textsl{Melee Weapon Attack:} 7 to hit, reach 5 ft., one target. \textsl{Hit:} 11 (2d6 + 4) slashing damage.

\end{DndMonster}



\begin{DndMonster}[float=t]{Water Weird}
\DndMonsterType{Large elemental, neutral}
\DndMonsterBasics[
armor-class = {13},
hit-points = {\DndDice{9d10 + 9}},
speed = {0 ft., swim 60 ft.}
]
\DndMonsterAbilityScores[
 str = 17,
 dex = 16,
 con = 13,
 int = 11,
 wis = 10,
 cha = 10
]
\DndMonsterDetails[
damage-resistances = {fire; bludgeoning, piercing, slashing from nonmagical attacks},
damage-immunities = {poison},
condition-immunities = {exhaustion, grappled, paralyzed, poisoned, restrained, prone, unconscious},
senses = {blindsight 30 ft., passive perception 10},
languages = {understands Aquan but doesn't speak},
challenge = 3,
]
\DndMonsterAction{Invisible in Water}
The water weird is invisible while fully immersed in water.

\DndMonsterAction{Water Bound}
The water weird dies if it leaves the water to which it is bound or if that water is destroyed.

\DndMonsterSection{Actions}
\DndMonsterAction{Constrict}
\textsl{Melee Weapon Attack:} 5 to hit, reach 10 ft., one creature. \textsl{Hit:} 13 (3d6 + 3) bludgeoning damage. If the target is Medium or smaller, it is grappled (escape DC 13) and pulled 5 feet toward the water weird. Until this grapple ends, the target is restrained, the water weird tries to drown it, and the water weird can't constrict another target.

\end{DndMonster}



\begin{DndMonster}[float=t]{Xorn}
\DndMonsterType{Medium elemental, neutral}
\DndMonsterBasics[
armor-class = {19 (natural armor)},
hit-points = {\DndDice{7d8 + 42}},
speed = {20 ft., burrow 20 ft.}
]
\DndMonsterAbilityScores[
 str = 17,
 dex = 10,
 con = 22,
 int = 11,
 wis = 10,
 cha = 11
]
\DndMonsterDetails[
skills = {Perception +6, Stealth +3},
damage-resistances = {piercing, slashing from nonmagical attacks that aren't adamantine},
senses = {darkvision 60 ft., tremorsense 60 ft., passive perception 16},
languages = {Terran},
challenge = 5,
]
\DndMonsterAction{Earth Glide}
The xorn can burrow through nonmagical, unworked earth and stone. While doing so, the xorn doesn't disturb the material it moves through.

\DndMonsterAction{Stone Camouflage}
The xorn has advantage on Dexterity (Stealth) checks made to hide in rocky terrain.

\DndMonsterAction{Treasure Sense}
The xorn can pinpoint, by scent, the location of precious metals and stones, such as coins and gems, within 60 feet of it.

\DndMonsterSection{Actions}
\DndMonsterAction{Multiattack}
The xorn makes three claw attacks and one bite attack.

\DndMonsterAction{Bite}
\textsl{Melee Weapon Attack:} 6 to hit, reach 5 ft., one target. \textsl{Hit:} 13 (3d6 + 3) piercing damage.

\DndMonsterAction{Claw}
\textsl{Melee Weapon Attack:} 6 to hit, reach 5 ft., one target. \textsl{Hit:} 6 (1d6 + 3) slashing damage.

\end{DndMonster}


\chapter{Spells}
\addcontentsline{toc}{section}{Antipathy/Sympathy}


\DndSpellHeader%
{Antipathy/Sympathy}
{8th level Enchantment}
{1 hour}
{60 feet}
{V, S, either a lump of alum soaked in vinegar for the antipathy effect or a drop of honey for the sympathy effect}
{10 days}
This spell attracts or repels creatures of your choice. You target something within range, either a Huge or smaller object or creature or an area that is no larger than a 200-foot cube. Then specify a kind of intelligent creature, such as red dragons, goblins, or vampires. You invest the target with an aura that either attracts or repels the specified creatures for the duration. Choose antipathy or sympathy as the aura's effect.

\subparagraph{Antipathy}
The enchantment causes creatures of the kind you designated to feel an intense urge to leave the area and avoid the target. When such a creature can see the target or comes within 60 feet of it, the creature must succeed on a Wisdom saving throw or become frightened. The creature remains frightened while it can see the target or is within 60 feet of it. While frightened by the target, the creature must use its movement to move to the nearest safe spot from which it can't see the target. If the creature moves more than 60 feet from the target and can't see it, the creature is no longer frightened, but the creature becomes frightened again if it regains sight of the target or moves within 60 feet of it.

\subparagraph{Sympathy}
The enchantment causes the specified creatures to feel an intense urge to approach the target while within 60 feet of it or able to see it. When such a creature can see the target or comes within 60 feet of it, the creature must succeed on a Wisdom saving throw or use its movement on each of its turns to enter the area or move within reach of the target. When the creature has done so, it can't willingly move away from the target.

If the target damages or otherwise harms an affected creature, the affected creature can make a Wisdom saving throw to end the effect, as described below.

\subparagraph{Ending the Effect}
If an affected creature ends its turn while not within 60 feet of the target or able to see it, the creature makes a Wisdom saving throw. On a successful save, the creature is no longer affected by the target and recognizes the feeling of repugnance or attraction as magical. In addition, a creature affected by the spell is allowed another Wisdom saving throw every 24 hours while the spell persists.

A creature that successfully saves against this effect is immune to it for 1 minute, after which time it can be affected again.

\addcontentsline{toc}{section}{Arms of Hadar}


\DndSpellHeader%
{Arms of Hadar}
{1st level Conjuration}
{1 action}
{Self (10 feet radius)}
{V, S}
{Instantaneous}
You invoke the power of Hadar, the Dark Hunger. Tendrils of dark energy erupt from you and batter all creatures within 10 feet of you. Each creature in that area must make a Strength saving throw. On a failed save, a target takes 2d6 necrotic damage and can't take reactions until its next turn. On a successful save, the creature takes half damage, but suffers no other effect.

\addcontentsline{toc}{section}{Banishment}


\DndSpellHeader%
{Banishment}
{4th level Abjuration}
{1 action}
{60 feet}
{V, S, an item distasteful to the target}
{Concentration, up to 1 minute}
You attempt to send one creature that you can see within range to another plane of existence. The target must succeed on a Charisma saving throw or be banished.

If the target is native to the plane of existence you're on, you banish the target to a harmless demiplane. While there, the target is incapacitated. The target remains there until the spell ends, at which point the target reappears in the space it left or in the nearest unoccupied space if that space is occupied.

If the target is native to a different plane of existence than the one you're on, the target is banished with a faint popping noise, returning to its home plane. If the spell ends before 1 minute has passed, the target reappears in the space it left or in the nearest unoccupied space if that space is occupied. Otherwise, the target doesn't return.

\addcontentsline{toc}{section}{Conjure Elemental}


\DndSpellHeader%
{Conjure Elemental}
{5th level Conjuration}
{1 minute}
{90 feet}
{V, S, burning incense for air, soft clay for earth, sulfur and phosphorus for fire, or water and sand for water}
{Concentration, up to 1 hour}
You call forth an elemental servant. Choose an area of air, earth, fire, or water that fills a 10-foot cube within range. An elemental of challenge rating 5 or lower appropriate to the area you chose appears in an unoccupied space within 10 feet of it. For example, a \textbf{fire elemental} emerges from a bonfire, and an \textbf{earth elemental} rises up from the ground. The elemental disappears when it drops to 0 hit points or when the spell ends.

The elemental is friendly to you and your companions for the duration. Roll initiative for the elemental, which has its own turns. It obeys any verbal commands that you issue to it (no action required by you). If you don't issue any commands to the elemental, it defends itself from hostile creatures but otherwise takes no actions.

If your concentration is broken, the elemental doesn't disappear. Instead, you lose control of the elemental, it becomes hostile toward you and your companions, and it might attack. An uncontrolled elemental can't be dismissed by you, and it disappears 1 hour after you summoned it.

The DM has the elemental's statistics.

\addcontentsline{toc}{section}{Control Water}


\DndSpellHeader%
{Control Water}
{4th level Transmutation}
{1 action}
{300 feet}
{V, S, a drop of water and a pinch of dust}
{Concentration, up to 10 minutes}
Until the spell ends, you control any freestanding water inside an area you choose that is a cube up to 100 feet on a side. You can choose from any of the following effects when you cast this spell. As an action on your turn, you can repeat the same effect or choose a different one.

\subparagraph{Flood}
You cause the water level of all standing water in the area to rise by as much as 20 feet. If the area includes a shore, the flooding water spills over onto dry land.

If you choose an area in a large body of water, you instead create a 20-foot tall wave that travels from one side of the area to the other and then crashes down. Any Huge or smaller vehicles in the wave's path are carried with it to the other side. Any Huge or smaller vehicles struck by the wave have a No effect chance of capsizing.

The water level remains elevated until the spell ends or you choose a different effect. If this effect produced a wave, the wave repeats on the start of your next turn while the flood effect lasts.

\subparagraph{Part Water}
You cause water in the area to move apart and create a trench. The trench extends across the spell's area, and the separated water forms a wall to either side. The trench remains until the spell ends or you choose a different effect. The water then slowly fills in the trench over the course of the next round until the normal water level is restored.

\subparagraph{Redirect Flow}
You cause flowing water in the area to move in a direction you choose, even if the water has to flow over obstacles, up walls, or in other unlikely directions. The water in the area moves as you direct it, but once it moves beyond the spell's area, it resumes its flow based on the terrain conditions. The water continues to move in the direction you chose until the spell ends or you choose a different effect.

\subparagraph{Whirlpool}
This effect requires a body of water at least 50 feet square and 25 feet deep. You cause a whirlpool to form in the center of the area. The whirlpool forms a vortex that is 5 feet wide at the base, up to 50 feet wide at the top, and 25 feet tall. Any creature or object in the water and within 25 feet of the vortex is pulled 10 feet toward it. A creature can swim away from the vortex by making a Strength (Athletics) check against your spell save DC.

When a creature enters the vortex for the first time on a turn or starts its turn there, it must make a Strength saving throw. On a failed save, the creature takes 2d8 bludgeoning damage and is caught in the vortex until the spell ends. On a successful save, the creature takes half damage, and isn't caught in the vortex. A creature caught in the vortex can use its action to try to swim away from the vortex as described above, but has disadvantage on the Strength (Athletics) check to do so.

The first time each turn that an object enters the vortex, the object takes 2d8 bludgeoning damage; this damage occurs each round it remains in the vortex.

\addcontentsline{toc}{section}{Create or Destroy Water}


\DndSpellHeader%
{Create or Destroy Water}
{1st level Transmutation}
{1 action}
{30 feet}
{V, S, a drop of water if creating water or a few grains of sand if destroying it}
{Instantaneous}
You either create or destroy water.

\subparagraph{Create Water}
You create up to 10 gallons of clean water within range in an open container. Alternatively, the water falls as rain in a 30-foot cube within range, extinguishing exposed flames in the area.

\subparagraph{Destroy Water}
You destroy up to 10 gallons of water in an open container within range. Alternatively, you destroy fog in a 30-foot cube within range.

\addcontentsline{toc}{section}{Detect Evil and Good}


\DndSpellHeader%
{Detect Evil and Good}
{1st level Divination}
{1 action}
{Self}
{V, S}
{Concentration, up to 10 minutes}
For the duration, you know if there is an aberration, celestial, elemental, fey, fiend, or undead within 30 feet of you, as well as where the creature is located. Similarly, you know if there is a place or object within 30 feet of you that has been magically consecrated or desecrated.

The spell can penetrate most barriers, but it is blocked by 1 foot of stone, 1 inch of common metal, a thin sheet of lead, or 3 feet of wood or dirt.

\addcontentsline{toc}{section}{Detect Magic}


\DndSpellHeader%
{Detect Magic}
{1st level Divination}
{1 action}
{Self}
{V, S}
{Concentration, up to 10 minutes}
For the duration, you sense the presence of magic within 30 feet of you. If you sense magic in this way, you can use your action to see a faint aura around any visible creature or object in the area that bears magic, and you learn its \textit{school of magic}, if any.

The spell can penetrate most barriers, but it is blocked by 1 foot of stone, 1 inch of common metal, a thin sheet of lead, or 3 feet of wood or dirt.

\addcontentsline{toc}{section}{Dispel Magic}


\DndSpellHeader%
{Dispel Magic}
{3rd level Abjuration}
{1 action}
{120 feet}
{V, S}
{Instantaneous}
Choose one creature, object, or magical effect within range. Any spell of 3rd level or lower on the target ends. For each spell of 4th level or higher on the target, make an ability check using your spellcasting ability. The DC equals 10 + the spell's level. On a successful check, the spell ends.

\addcontentsline{toc}{section}{Fog Cloud}


\DndSpellHeader%
{Fog Cloud}
{1st level Conjuration}
{1 action}
{120 feet}
{V, S}
{Concentration, up to 1 hour}
You create a 20-foot-radius sphere of fog centered on a point within range. The sphere spreads around corners, and its area is Vision and Light. It lasts for the duration or until a wind of moderate or greater speed (at least 10 miles per hour) disperses it.

\addcontentsline{toc}{section}{Gaseous Form}


\DndSpellHeader%
{Gaseous Form}
{3rd level Transmutation}
{1 action}
{Touch}
{V, S, a bit of gauze and a wisp of smoke}
{Concentration, up to 1 hour}
You transform a willing creature you touch, along with everything it's wearing and carrying, into a misty cloud for the duration. The spell ends if the creature drops to 0 hit points. An incorporeal creature isn't affected.

While in this form, the target's only method of movement is a flying speed of 10 feet. The target can enter and occupy the space of another creature. The target has resistance to nonmagical damage, and it has advantage on Strength, Dexterity, and Constitution saving throws. The target can pass through small holes, narrow openings, and even mere cracks, though it treats liquids as though they were solid surfaces. The target can't fall and remains hovering in the air even when stunned or otherwise incapacitated.

While in the form of a misty cloud, the target can't talk or manipulate objects, and any objects it was carrying or holding can't be dropped, used, or otherwise interacted with. The target can't attack or cast spells.

\addcontentsline{toc}{section}{Greater Restoration}


\DndSpellHeader%
{Greater Restoration}
{5th level Abjuration}
{1 action}
{Touch}
{V, S, diamond dust worth at least 100 gp, which the spell consumes}
{Instantaneous}
You imbue a creature you touch with positive energy to undo a debilitating effect. You can reduce the target's exhaustion level by one, or end one of the following effects on the target:


\begin{itemize}
\item One effect that charmed or petrified the target
\item One curse, including the target's attunement to a cursed magic item
\item Any reduction to one of the target's ability scores
\item One effect reducing the target's hit point maximum
\end{itemize}

\addcontentsline{toc}{section}{Gust of Wind}


\DndSpellHeader%
{Gust of Wind}
{2nd level Evocation}
{1 action}
{Self (60 feet line)}
{V, S, a legume seed}
{Concentration, up to 1 minute}
A line of strong wind 60 feet long and 10 feet wide blasts from you in a direction you choose for the spell's duration. Each creature that starts its turn in the line must succeed on a Strength saving throw or be pushed 15 feet away from you in a direction following the line.

Any creature in the line must spend 2 feet of movement for every 1 foot it moves when moving closer to you.

The gust disperses gas or vapor, and it extinguishes candles, torches, and similar unprotected flames in the area. It causes protected flames, such as those of lanterns, to dance wildly and has a No effect chance to extinguish them.

As a bonus action on each of your turns before the spell ends, you can change the direction in which the line blasts from you.

\addcontentsline{toc}{section}{Identify}


\DndSpellHeader%
{Identify}
{1st level Divination}
{1 minute}
{Touch}
{V, S, a pearl worth at least 100 gp and an owl feather}
{Instantaneous}
You choose one object that you must touch throughout the casting of the spell. If it is a magic item or some other magic-imbued object, you learn its properties and how to use them, whether it requires attunement to use, and how many charges it has, if any. You learn whether any spells are affecting the item and what they are. If the item was created by a spell, you learn which spell created it.

If you instead touch a creature throughout the casting, you learn what spells, if any, are currently affecting it.

\addcontentsline{toc}{section}{Invisibility}


\DndSpellHeader%
{Invisibility}
{2nd level Illusion}
{1 action}
{Touch}
{V, S, an eyelash encased in gum arabic}
{Concentration, up to 1 hour}
A creature you touch becomes invisible until the spell ends. Anything the target is wearing or carrying is invisible as long as it is on the target's person. The spell ends for a target that attacks or casts a spell.

\addcontentsline{toc}{section}{Levitate}


\DndSpellHeader%
{Levitate}
{2nd level Transmutation}
{1 action}
{60 feet}
{V, S, either a small leather loop or a piece of golden wire bent into a cup shape with a long shank on one end}
{Concentration, up to 10 minutes}
One creature or loose object of your choice that you can see within range rises vertically, up to 20 feet, and remains suspended there for the duration. The spell can levitate a target that weighs up to 500 pounds. An unwilling creature that succeeds on a Constitution saving throw is unaffected.

The target can move only by pushing or pulling against a fixed object or surface within reach (such as a wall or a ceiling), which allows it to move as if it were climbing. You can change the target's altitude by up to 20 feet in either direction on your turn. If you are the target, you can move up or down as part of your move. Otherwise, you can use your action to move the target, which must remain within the spell's range.

When the spell ends, the target floats gently to the ground if it is still aloft.

\addcontentsline{toc}{section}{Move Earth}


\DndSpellHeader%
{Move Earth}
{6th level Transmutation}
{1 action}
{120 feet}
{V, S, an iron blade and a small bag containing a mixture of soils—clay, loam, and sand}
{Concentration, up to 2 hours}
Choose an area of terrain no larger than 40 feet on a side within range. You can reshape dirt, sand, or clay in the area in any manner you choose for the duration. You can raise or lower the area's elevation, create or fill in a trench, erect or flatten a wall, or form a pillar. The extent of any such changes can't exceed half the area's largest dimension. So, if you affect a 40-foot square, you can create a pillar up to 20 feet high, raise or lower the square's elevation by up to 20 feet, dig a trench up to 20 feet deep, and so on. It takes 10 minutes for these changes to complete.

At the end of every 10 minutes you spend concentration on the spell, you can choose a new area of terrain to affect.

Because the terrain's transformation occurs slowly, creatures in the area can't usually be trapped or injured by the ground's movement.

This spell can't manipulate natural stone or stone construction. Rocks and structures shift to accommodate the new terrain. If the way you shape the terrain would make a structure unstable, it might collapse.

Similarly, this spell doesn't directly affect plant growth. The moved earth carries any plants along with it.

\addcontentsline{toc}{section}{Passwall}


\DndSpellHeader%
{Passwall}
{5th level Transmutation}
{1 action}
{30 feet}
{V, S, a pinch of sesame seeds}
{1 hour}
A passage appears at a point of your choice that you can see on a wooden, plaster, or stone surface (such as a wall, a ceiling, or a floor) within range, and lasts for the duration. You choose the opening's dimensions: up to 5 feet wide, 8 feet tall, and 20 feet deep. The passage creates no instability in a structure surrounding it.

When the opening disappears, any creatures or objects still in the passage created by the spell are safely ejected to an unoccupied space nearest to the surface on which you cast the spell.

\addcontentsline{toc}{section}{Phantasmal Killer}


\DndSpellHeader%
{Phantasmal Killer}
{4th level Illusion}
{1 action}
{120 feet}
{V, S}
{Concentration, up to 1 minute}
You tap into the nightmares of a creature you can see within range and create an illusory manifestation of its deepest fears, visible only to that creature. The target must make a Wisdom saving throw. On a failed save, the target becomes frightened for the duration. At the end of each of the target's turns before the spell ends, the target must succeed on a Wisdom saving throw or take 4d10 psychic damage. On a successful save, the spell ends.

\addcontentsline{toc}{section}{Plane Shift}


\DndSpellHeader%
{Plane Shift}
{7th level Conjuration}
{1 action}
{Touch}
{V, S, a forked, metal rod worth at least 250 gp, attuned to a particular plane of existence}
{Instantaneous}
You and up to eight willing creatures who link hands in a circle are transported to a different plane of existence. You can specify a target destination in general terms, such as the City of Brass on the Elemental Plane of Fire or the palace of Dispater on the second level of the Nine Hells, and you appear in or near that destination. If you are trying to reach the City of Brass, for example, you might arrive in its Street of Steel, before its Gate of Ashes, or looking at the city from across the Sea of Fire, at the DM's discretion.

Alternatively, if you know the sigil sequence of a teleportation circle on another plane of existence, this spell can take you to that circle. If the teleportation circle is too small to hold all the creatures you transported, they appear in the closest unoccupied spaces next to the circle.

You can use this spell to banish an unwilling creature to another plane. Choose a creature within your reach and make a melee spell attack against it. On a hit, the creature must make a Charisma saving throw. If the creature fails this save, it is transported to a random location on the plane of existence you specify. A creature so transported must find its own way back to your current plane of existence.

\addcontentsline{toc}{section}{Purify Food and Drink}


\DndSpellHeader%
{Purify Food and Drink}
{1st level Transmutation}
{1 action}
{10 feet}
{V, S}
{Instantaneous}
All nonmagical food and drink within a 5-foot-radius sphere centered on a point of your choice within range is purified and rendered free of poison and disease.

\addcontentsline{toc}{section}{Stone Shape}


\DndSpellHeader%
{Stone Shape}
{4th level Transmutation}
{1 action}
{Touch}
{V, S, soft clay, which must be worked into roughly the desired shape of the stone object}
{Instantaneous}
You touch a stone object of Medium size or smaller or a section of stone no more than 5 feet in any dimension and form it into any shape that suits your purpose. So, for example, you could shape a large rock into a weapon, idol, or coffer, or make a small passage through a wall, as long as the wall is less than 5 feet thick. You could also shape a stone door or its frame to seal the door shut. The object you create can have up to two hinges and a latch, but finer mechanical detail isn't possible.

\addcontentsline{toc}{section}{Tongues}


\DndSpellHeader%
{Tongues}
{3rd level Divination}
{1 action}
{Touch}
{V, a small clay model of a ziggurat}
{1 hour}
This spell grants the creature you touch the ability to understand any spoken language it hears. Moreover, when the target speaks, any creature that knows at least one language and can hear the target understands what it says.

\addcontentsline{toc}{section}{Wall of Stone}


\DndSpellHeader%
{Wall of Stone}
{5th level Evocation}
{1 action}
{120 feet}
{V, S, a small block of granite}
{Concentration, up to 10 minutes}
A nonmagical wall of solid stone springs into existence at a point you choose within range. The wall is 6 inches thick and is composed of ten 10-foot-by-10-foot panels. Each panel must be contiguous with at least one other panel. Alternatively, you can create 10-foot-by-20-foot panels that are only 3 inches thick.

If the wall cuts through a creature's space when it appears, the creature is pushed to one side of the wall (your choice). If a creature would be surrounded on all sides by the wall (or the wall and another solid surface), that creature can make a Dexterity saving throw. On a success, it can use its reaction to move up to its speed so that it is no longer enclosed by the wall.

The wall can have any shape you desire, though it can't occupy the same space as a creature or object. The wall doesn't need to be vertical or rest on any firm foundation. It must, however, merge with and be solidly supported by existing stone. Thus, you can use this spell to bridge a chasm or create a ramp.

If you create a span greater than 20 feet in length, you must halve the size of each panel to create supports. You can crudely shape the wall to create crenellations, battlements, and so on.

The wall is an object made of stone that can be damaged and thus breached. Each panel has AC 15 and 30 hit points per inch of thickness. Reducing a panel to 0 hit points destroys it and might cause connected panels to collapse at the DM's discretion.

If you maintain your concentration on this spell for its whole duration, the wall becomes permanent and can't be dispelled. Otherwise, the wall disappears when the spell ends.

\addcontentsline{toc}{section}{Water Breathing}


\DndSpellHeader%
{Water Breathing}
{3rd level Transmutation}
{1 action}
{30 feet}
{V, S, a short reed or piece of straw}
{24 hours}
This spell grants up to ten willing creatures you can see within range the ability to breathe underwater until the spell ends. Affected creatures also retain their normal mode of respiration.

\addcontentsline{toc}{section}{Water Walk}


\DndSpellHeader%
{Water Walk}
{3rd level Transmutation}
{1 action}
{30 feet}
{V, S, a piece of cork}
{1 hour}
This spell grants the ability to move across any liquid surface—such as water, acid, mud, snow, quicksand, or lava—as if it were harmless solid ground (creatures crossing molten lava can still take damage from the heat). Up to ten willing creatures you can see within range gain this ability for the duration.

If you target a creature submerged in a liquid, the spell carries the target to the surface of the liquid at a rate of 60 feet per round.

\addcontentsline{toc}{section}{Web}


\DndSpellHeader%
{Web}
{2nd level Conjuration}
{1 action}
{60 feet}
{V, S, a bit of spiderweb}
{Concentration, up to 1 hour}
You conjure a mass of thick, sticky webbing at a point of your choice within range. The webs fill a 20-foot cube from that point for the duration. The webs are difficult terrain and lightly obscure their area.

If the webs aren't anchored between two solid masses (such as walls or trees) or layered across a floor, wall, or ceiling, the conjured web collapses on itself, and the spell ends at the start of your next turn. Webs layered over a flat surface have a depth of 5 feet.

Each creature that starts its turn in the webs or that enters them during its turn must make a Dexterity saving throw. On a failed save, the creature is restrained as long as it remains in the webs or until it breaks free.

A creature restrained by the webs can use its action to make a Strength check against your spell save DC. If it succeeds, it is no longer restrained.

The webs are flammable. Any 5-foot cube of webs exposed to fire burns away in 1 round, dealing 2d4 fire damage to any creature that starts its turn in the fire.

\end{document}
